\subsection{The measure of degenerate flag fibres}
\label{shared-degenerate-flag-fibres}

We now supply the geometric change of variables needed between
the original flag integral and the already proved Euclidean
small Gram determinant estimate~\cite{global-gram-degenerate-euclidean-volume}.
The stabilizer measure remains
explicit.  In particular, no probability measure on a noncompact
unipotent group is introduced.  The construction and normalization
are those of the pinned source~\cite{global-gram-degenerate-fibres-source}.

\paragraph{The original coordinates.}
Put $V=E_0\oplus W\oplus F_0$, where $E_0,F_0$ have bases
$e_1,\ldots,e_r$ and $f_1,\ldots,f_r$, and
$W=\mathbb C^{k+1}$.  All these coordinates are orthonormal for
the Euclidean norm.  The indefinite Hermitian form, linear in its
first argument, has matrix
\[
 J=\begin{pmatrix}0&0&I_r\\0&J_W&0\\I_r&0&0\end{pmatrix},
 \qquad J_W=\operatorname{diag}(1,\ldots,1,-1).
\]
Write $h(v,w)=w^*Jv$, $G=U(J)$ and $K=G\cap U(V)$.
The original dual flag coordinates are here listed in matching
order.  Reversing their printed order is an orthonormal coordinate
permutation, with real Jacobian one.

For an upper unitriangular matrix $u$, a matrix
$B=(B_1,\ldots,B_r)\in\operatorname{Hom}(\mathbb C^r,W)$,
and a skew Hermitian matrix $S$, set
\begin{align*}
 U(u)&=\operatorname{diag}(u,I_{k+1},u^{-*}),\\
 P(B,S)&=\begin{pmatrix}
 I_r&-B^*J_W&S-\tfrac12B^*J_WB\\
 0&I_{k+1}&B\\0&0&I_r
 \end{pmatrix}.
\end{align*}
The group $N$ consists of the products $U(u)P(B,S)$.
Its original Haar measure $dn$ is product Lebesgue measure in
the complex entries $u_{ab}$ for $a<b$, the complex coordinates
of $B$, the complex entries $S_{ab}$ for $a<b$, and the real
entries $-iS_{aa}$.  These formulas specify the measure, including
all factors of two.

For a negative vector $z\in W$, let
$v_1=e_1,\ldots,v_r=e_r,v_{r+1}=z$ and define
\begin{align*}
 L_i(z,x)&=\operatorname{dist}\bigl(x^{-1}v_i,
       x^{-1}\operatorname{span}(v_1,\ldots,v_{i-1})\bigr),\\
 S_z(x)&=\max\bigl(1,L_i(z,x),L_i(z,x)^{-1}:1\le i\le r+1\bigr).
\end{align*}
All these lengths are positive.  Left multiplication by $N$ and
right multiplication by $K$ preserve them: the former changes a
flag vector only by earlier vectors, and the latter is Euclidean
unitary.  If
\[
 t=\operatorname{diag}(a_1,\ldots,a_r,m,a_1^{-1},\ldots,a_r^{-1}),
 \qquad a_i>0,\quad m\in U(J_W),
\]
then $L_i(z,nt)=a_i^{-1}$ for $i\le r$ and $n\in N$.

\paragraph{The common normal form.}
Fix $1\le j\le r$ and $j+2\le k$.  Partition
$\{1,\ldots,j\}$ into consecutive blocks; if $j<r$, append
$\{j+1,\ldots,r\}$ as one further block.  Suppose $g_0\in G$
and negative $z,z'\in W$ satisfy the following original normal
form conditions.  There is $c\in\mathbb C$, $|c|=1$, such that
$g_0z'=cz$; both $g_0$ and $g_0^{-1}$ send
$E_j=\operatorname{span}(e_1,\ldots,e_j)$ into
$F_j=\operatorname{span}(f_1,\ldots,f_j)$; and
$g_0e_a=ce_a$, $g_0f_a=cf_a$ for $a>j$.
Finally, the $j$ by $j$ matrix
\[
 Z_{ab}=h(g_0e_b,e_a)
\]
is zero outside its diagonal blocks, and within each block it is
antidiagonal with nonzero entries of a common absolute value.
The phases of those entries need not agree.  Let $w$ reverse
each block and write $\zeta_a=Z_{a,w(a)}$.

Let $D_l(g)$ be the determinant of the leading $l$ by $l$
submatrix of $(h(ge_b,e_a))_{a,b\le r}$, with $D_0=1$.
Put $D=|D_j(g_0)|>0$, and suppose $\Delta\ge1$ satisfies
\[
 D_l(g_0)\ne0\quad\Longrightarrow\quad
              |D_l(g_0)|^{-1}\le\Delta\qquad(0\le l\le r).
\]
These are precisely the nonzero leading-minor bounds used below.

The following measure and functions also specify the conjugated
normal form.  Take $\gamma=n_1g_0n_2$ with $n_1,n_2\in N$.
Let $R_\gamma=N\cap\gamma^{-1}N\gamma$.  At $g_0$, write
$u_{\rm in}$ for an upper unitriangular matrix supported inside
the blocks, and let $B_{\rm top},S_{\rm top}$ be supported on
indices greater than $j$.  Give
\[
 R_0=\{U(u_{\rm in})P(B_{\rm top},S_{\rm top})\}
\]
product coordinate Lebesgue measure $d\rho$.  The prescribed
measure on $R_\gamma$ is the image of $s\,d\rho$, $s>0$,
under $\rho\mapsto n_2^{-1}\rho n_2$.
We prove below that $R_0=N\cap g_0^{-1}Ng_0$, so this is
the original common-stabilizer measure.

Write $\mathcal C_\gamma$ for the common $h$-orthogonal
complement of $E_0+\mathbb Cz'$ and
$\gamma^{-1}(E_0+\mathbb Cz)$.  Let $Q_\gamma(x)$ be the
Euclidean orthogonal projection onto
\[
       \ker(J-I)\cap x^{-1}(\mathcal C_\gamma+E_0).
\]
For vectors $b_1,\ldots,b_j$ and a Euclidean orthogonal
projection $Q$, put
$\operatorname{Gram}_Q(b_1,\ldots,b_j)
 =\|Qb_1\wedge\cdots\wedge Qb_j\|$.  Define the original
normalized Gram volume by
\[
 Y_\gamma(x)=\frac{
    \operatorname{Gram}_{Q_\gamma(x)}
       (x^{-1}\gamma^{-1}e_1,\ldots,x^{-1}\gamma^{-1}e_j)}D.
\]
Unipotent multiplication preserves the leading minors, so the
same $D$ and $\Delta$ apply to $\gamma$.

\begin{theorem}[The original degenerate fibre estimate]
\label{thm:original-degenerate-fibre-estimate}
Retain all the hypotheses and measures just specified.  Let
$\varphi:G\to[0,\infty]$ be measurable and satisfy
\[
       \int_{R_\gamma}\varphi(\rho x)\,d\mu_R(\rho)
                           \le1\qquad(x\in G).
\]
For every $g\in G$, the set
$\{n\in N:Y_\gamma(ng)=0\}$ has $dn$-measure zero.
For $R\ge1$ and $y>0$, one has
\begin{align}
 &\int_N\varphi(ng)
    {\bf1}_{\{S_{z'}(ng),S_z(\gamma ng)\le R,\;
                         y<Y_\gamma(ng)\le2y\}}\,dn \notag\\
 &\quad\le\frac{(4r)^{4r^2}}s
       (R^{2r}\Delta)^{8r^2}
       \left(\frac{\pi e}{k-1-j}\right)^{j(k-1-j)}
                         (2y)^{2(k-1-j)}.
 \label{eq:original-degenerate-fibre-estimate}
\end{align}
\end{theorem}

\begin{proof}
We first take $\gamma=g_0$ and $g=t$ of the displayed Levi
form, and then prove the transport to arbitrary $\gamma,g$.

\emph{The quotient coordinates have Jacobian one.}
Let $u_{\rm off}$ be upper unitriangular, with zero entries
inside every diagonal block.  Let $B_{\rm sec}$ have zero
columns with index greater than $j$, and let $S_{\rm sec}$ have zero entries
when both indices exceed $j$.  The section is
\[
           \sigma=P(B_{\rm sec},S_{\rm sec})U(u_{\rm off}).
\]
Every element of $N$ has a unique expression $\rho\sigma$,
with $\rho\in R_0$, and
\begin{equation}\label{eq:degenerate-measure-split}
                         dn=d\rho\,d\sigma,
\end{equation}
where $d\sigma$ is the product of the remaining coordinate
Lebesgue measures.  Here is a verification of the Jacobian.
An upper unitriangular $u$ factors uniquely as
$u=u_{\rm in}u_{\rm off}$: its diagonal blocks determine
$u_{\rm in}$, and each off-diagonal block is multiplied by
an inverse unitriangular diagonal block.  All the complex
determinants of these linear changes are one.
Moreover
\begin{align*}
 P(B,S)P(C,T)
   &=P\bigl(B+C,S+T+\tfrac12(C^*J_WB-B^*J_WC)\bigr),\\
 U(u)P(B,S)U(u)^{-1}&=P(Bu^*,uSu^*).
\end{align*}
The first identity splits the top and section variables by
translations in the skew coordinates, so its real Jacobian is
one.  In the second identity $u$ is unitriangular.  Its action
on $B$ has determinant one, and its action on the real space
of skew Hermitian matrices is unitriangular after ordering the
coordinate pairs from the last indices to the first.  Thus
that real determinant is also one.  Moving $u_{\rm in}$ past
$P(B,S)$ and using these two identities proves
\eqref{eq:degenerate-measure-split}.  Left and right
multiplication in upper unitriangular coordinates are triangular
affine changes with diagonal one.  Together with the displayed
radical identities, this also proves directly that $dn$ and
$d\rho$ are left and right Haar measures.

For completeness, the stabilizer assertion follows from the
same matrices.  The normal form gives
\begin{equation}\label{eq:normal-moved-dual}
                     g_0^{-1}e_a=\overline{\zeta_a}f_{w(a)}
                             \qquad(a\le j).
\end{equation}
If $U(u)P(B,S)$ belongs to $g_0^{-1}Ng_0$, its action on
these vectors is upper unitriangular in their displayed order.
The middle and $E_0$ coordinates of its $f_b$ column, $b\le j$,
therefore vanish.  The block formulas give $B_b=0$ and
$S_{ab}=0$ for every $a$ and $b\le j$.  Its remaining
$f$ coordinates are $\overline{(u^{-1})_{ba}}$.  Since $w$
reverses the order within each block and preserves the order
of the blocks, the two triangularity conditions force
$u^{-1}$, and hence $u$, to be block diagonal.  These are
exactly the coordinates of $R_0$.  Conversely, a matrix with
these coordinates is sent into $N$ by conjugation with $g_0$:
within a block, reversal changes lower triangularity back to
upper triangularity; the remaining radical acts only on
$W$ and the last block.  Thus $R_0=N\cap g_0^{-1}Ng_0$.

\emph{Bounding the section coordinates from the two flags.}
Suppose $S_{z'}(\sigma t),S_z(g_0\sigma t)\le R$, and put
$M=R^{2r}\Delta\ge1$.  First note that
\begin{equation}\label{eq:normal-prefix-product}
                       \left(\prod_{a=1}^{l}|\zeta_a|\right)^{-1}
                              \le\Delta\qquad(0\le l\le j).
\end{equation}
At a block endpoint this product is the absolute value of the
corresponding nonzero leading minor.  Within a block its
logarithm interpolates linearly between the two endpoint
logarithms, since all its entries have the same absolute value.
The estimate follows from the endpoint bounds, including the
empty product one.

For $b\le j$ we claim
\begin{equation}\label{eq:normal-dual-pair-bound}
     |h(\sigma^{-1}f_b,v)|\le M
            \quad\hbox{whenever }\|t^{-1}v\|\le R.
\end{equation}
Write $b=w(l)$ and use the first $l$ vectors
$(g_0\sigma t)^{-1}e_a$.  Their exterior product has norm at
most $R^l$, since their Gram--Schmidt lengths are the first
$l$ original flag lengths.  Pair them with
$t^{-1}e_{w(1)},\ldots,t^{-1}e_{w(l-1)},t^{-1}v$.
Each test vector has norm at most $R$: for the first ones use
$a_i^{-1}=L_i(z',\sigma t)\le R$.  As $J$ is Euclidean
unitary, the absolute determinant of these Hermitian pairings
is at most $R^{2l}$ by the exterior Cauchy--Schwarz inequality.
By \eqref{eq:normal-moved-dual}, its first $l-1$ rows are
lower triangular with diagonal entries
$\overline{\zeta_1},\ldots,\overline{\zeta_{l-1}}$.
Indeed
$h(\sigma^{-1}f_b,e_a)=\overline{(u_{\rm off})_{ba}}$;
entries above this diagonal vanish either by upper
triangularity or by the zero entries within each block.
The determinant consequently has absolute value
\[
       \left(\prod_{a=1}^l|\zeta_a|\right)
                           |h(\sigma^{-1}f_b,v)|.
\]
Equation~\eqref{eq:normal-prefix-product} and $l\le r$ prove
\eqref{eq:normal-dual-pair-bound}.

Write $u=u_{\rm off}$, $B=B_{\rm sec}$, $S=S_{\rm sec}$.
The three components of $\sigma^{-1}f_b$ are
\[
 \left(u^{-1}\left(-S_{\cdot b}-\tfrac12(B^*J_WB)_{\cdot b}\right),
                      -B_b,\;u^*f_b\right).
\]
Taking $v=e_a$ in \eqref{eq:normal-dual-pair-bound} bounds
$|u_{ba}|$ by $M$ for $b\le j$.  The later rows of $u$
are identity rows, so every entry of $u$ has absolute value
at most $M$.  Taking $v=mw\in W$ with $\|w\|\le1$ gives
\[
                             \|m^{-1}B_b\|\le M.
\]
Here $\sup_{\|w\|\le1}|h_W(v,w)|=\|v\|$ because $J_W$
is Euclidean unitary.  Finally take $v=f_a$, using
$a_a=L_a(z',\sigma t)^{-1}\le R$.  Every coordinate of
$u^{-1}(-S_{\cdot b}-\tfrac12(B^*J_WB)_{\cdot b})$ is at
most $M$ in absolute value.  Multiplying by $u$ bounds its
untransformed coordinates by $rM^2$.  Since $m\in U(J_W)$,
\[
       |(B^*J_WB)_{ab}|\le\|m^{-1}B_a\|\|m^{-1}B_b\|\le M^2.
\]
Therefore $|S_{ab}|\le(r+\tfrac12)M^2\le2rM^2$ whenever
one index is at most $j$.  We have proved the precise boxes
\begin{equation}\label{eq:degenerate-coordinate-boxes}
 |u_{ab}|\le M,\qquad \|m^{-1}B_a\|\le M,\qquad
 |S_{ab}|,\;|-iS_{aa}|\le2rM^2
\end{equation}
for their respective section coordinates.

\emph{The projected Gram volume has no further Jacobian.}
The normal form implies
\[
 \mathcal C_{g_0}+E_0=(E_0+\mathbb Cz')^{\perp_h}.
\]
Indeed $g_0^{-1}E_j=F_j$, it preserves the remaining
isotropic lines, and $g_0^{-1}z$ is a nonzero multiple of
$z'$.  Adding $E_0$ removes exactly the extra orthogonality
conditions imposed by $F_j$.
If $v$ lies in $\ker(J-I)$ and
$\sigma tv\in\mathcal C_{g_0}+E_0$, its $F_0$ coordinates
are zero; its $E_0$ coordinates are then also zero because
$Jv=v$.  Thus $v\in W$, its last coordinate is zero, and
$h_W(mv,z')=0$.  Conversely these conditions imply membership.
Consequently $Q_{g_0}(\sigma t)$ is the projection in $W$ onto
\[
       W_0=\bigl(\operatorname{span}_{\mathbb C}
             \{e_{k+1},m^*J_Wz'\}\bigr)^\perp,
          \qquad \dim_{\mathbb C}W_0\ge k-1.
\]
It is independent of the section coordinates.  The middle
component of $(\sigma t)^{-1}f_b$ is $-m^{-1}B_b$.
Equation~\eqref{eq:normal-moved-dual} shows that the
determinant of the change of the first $j$ moved vectors is
$D$.  It cancels the denominator in $Y$.  Hence, writing
$Q_0$ for the orthogonal projection onto $W_0$,
\begin{equation}\label{eq:actual-degenerate-gram-coordinate}
       Y_{g_0}(\sigma t)=
              \operatorname{Gram}_{Q_0}(m^{-1}B_1,\ldots,m^{-1}B_j).
\end{equation}
Taking determinants of $m^*J_Wm=J_W$ gives $|\det m|=1$.
The real Jacobian of $B_a\mapsto m^{-1}B_a$ is therefore
exactly one, even when $m$ has large Euclidean norm.

\emph{Integration of the boxes.}
The functions $S_{z'}(x)$, $S_z(g_0x)$ and $Y_{g_0}(x)$
are invariant under left multiplication by $R_0$.
For $Y$, this follows because $R_0$ preserves both flags,
acts upper unitriangularly on their first $j$ vectors, and
preserves $\mathcal C_{g_0}+E_0$.  Therefore
\eqref{eq:degenerate-measure-split} and the assumption on
$\varphi$ reduce the integral to $1/s$ times the measure
of the allowable section coordinates.
By (GV8) and \eqref{eq:actual-degenerate-gram-coordinate}, the
volume of their $B$ coordinates is at most
\[
 \frac{\pi^{j^2+2j}M^{2j^2+4j}}{((j+2)!)^j}
       \left(\frac{\pi e}{k-1-j}\right)^{j(k-1-j)}
                             (2y)^{2(k-1-j)}.
\]
There are $p\le r(r-1)/2$ complex skew coordinates,
$j\le r$ real skew coordinates, and
$h\le r(r-1)/2$ complex upper coordinates.  Their box
volumes multiply the preceding expression by
\[
             (4\pi r^2M^4)^p(4rM^2)^j(\pi M^2)^h.
\]
The exponent of $M$ is at most $5r^2+3r\le8r^2$.
The exponent of $\pi$ is at most $2r^2+r\le3r^2$,
and the remaining factors $r$ and $2$ are bounded by
$(4r)^{2p+j}\le(4r)^{r^2}$.  Using $\pi<4$ and dropping
the factorial denominator proves the required factor
$(4r)^{4r^2}M^{8r^2}$.  The zero Gram locus has measure
zero by the same Euclidean Gram theorem and Fubini; all
coordinate measures here are sigma-finite.

\emph{Transport of the original measure.}
For arbitrary $g\in G$, write
$g=n_0t\kappa$ with $n_0\in N$, $\kappa\in K$ and $t$
as above.  Here is the elementary flag Iwasawa argument.
Euclidean Gram--Schmidt gives an orthonormal frame adapted to
$g^{-1}E_1\subset\cdots\subset g^{-1}E_r$.  Its projections
to the positive and negative eigenspaces of $J$, multiplied by
$\sqrt2$, are orthonormal frames: all Hermitian pairings in
the isotropic flag vanish.  Independent unitaries of these
two eigenspaces therefore give an element of $K$ carrying
the standard isotropic frame to this frame.  Adjust its
individual phases so that the diagonal entries of the
resulting action on $E_0$ are positive.  The transformed
matrix preserves the standard flag.  Its action on $E_0$
has positive diagonal $a$, and its induced action on
$E_0^{\perp_h}/E_0$ is $m\in U(J_W)$.  Dividing by $t$
leaves precisely a matrix of $N$.
The change $n\mapsto n_2nn_0$ preserves $dn$, since $N$
is unimodular.  Right multiplication by $K$ preserves the
lengths and the projected Gram volume; $K$ commutes with $J$.
Left multiplication by $n_1,n_2$ changes each leading flag
matrix by unitriangular row or column operations.  It thus
preserves every $D_l$, and
\[
       Y_{n_1g_0n_2}(x)=Y_{g_0}(n_2x).
\]
More explicitly,
$\mathcal C_{n_1g_0n_2}+E_0=n_2^{-1}(\mathcal C_{g_0}+E_0)$,
and the first $j$ moved vectors differ by a unitriangular
matrix; this proves the last identity including its
normalization.
Finally set
$\psi(x)=s\,\varphi(n_2^{-1}x\kappa)$ for $x\in G$.
The definition of $\mu_R$ gives
$\int_{R_0}\psi(\rho x)\,d\rho\le1$ for every $x\in G$.
Apply the already proved normal-form estimate to $\psi$, with
stabilizer scalar one, and
change variables by $n_2nn_0$.  Division by $s$ gives
\eqref{eq:original-degenerate-fibre-estimate}.  The same
measure-preserving change transports the null-set assertion.
\end{proof}

The theorem closes the original normal-form fibre supplier,
including the stabilizer scalar and arbitrary right translates.
Its hypotheses still require the original normal form and its
arithmetic measure.  The ensuing local orbital estimate is supplied by
Theorem~\ref{thm:actual-local-degenerate-bound}, using the original
flag-length tails.  The full arithmetic diagonal zero-Gram sum,
including the weighted count of leading minors, remains separate.
